\chapter*{ВВЕДЕНИЕ}
\addcontentsline{toc}{chapter}{ВВЕДЕНИЕ}
\sloppy

\textbf{Актуальность темы исследования.}
Большие языковые модели (LLM) всё активнее встраиваются в системы поддержки принятия решений: от медицинской диагностики и юридического анализа до финансового консультирования и автоматизированного обучения~\cite{bommasani2021opportunities}. Целевой аудиторией подобных систем являются разработчики ML-приложений, аналитики данных и исследователи, проектирующие инструменты на основе LLM для доменов с высокой ценой ошибки. По мере того как область применения подобных систем расширяется, принципиальное значение приобретает вопрос об их когнитивной надёжности: способны ли модели рассуждать нормативно корректно, или же они воспроизводят систематические ошибки, характерные для человеческого мышления? Когнитивные предубеждения~--- устойчивые отклонения от нормативных стандартов вероятностного и логического рассуждения~--- давно являются предметом исследований в психологии и поведенческой экономике~\cite{kahneman2011thinking}. С появлением масштабных языковых моделей, обученных на огромных корпусах текстов, порождённых людьми, закономерно возникает гипотеза о том, что подобные предубеждения могут быть имплицитно усвоены в ходе обучения и проявляться при решении задач, требующих строгого вероятностного или логического мышления. Результаты тестируемых моделей интерпретируются относительно трёх ориентиров: нормативного байесовского агента (верхняя граница), эмпирически измеренного уровня человека ($\approx$80\% ошибки конъюнкции~\cite{tversky1983extensional}, $\approx$10\% верных ответов в задаче Уэйсона~\cite{wason1966reasoning}) и случайного агента (нижняя граница). Работа носит междисциплинарный характер, объединяя методологию когнитивной психологии с эмпирическим исследованием систем искусственного интеллекта.

Несмотря на растущий массив литературы, посвящённой рациональности LLM~\cite{binz2023using, suri2024large, macmillan2024irrationality}, большинство существующих работ сосредоточено на отдельных задачах или узком наборе моделей, не рассматривая предубеждения как феномен, существующий на спектре~--- от вероятностного суждения к классическому логическому рассуждению и далее к индуктивному открытию правил. Настоящая работа восполняет этот пробел, предлагая систематическое сравнительное исследование трёх парадигм когнитивной психологии на девяти современных языковых моделях.

\textbf{Степень разработанности проблемы.}
Первые систематические исследования когнитивных предубеждений в LLM появились в 2023--2024 годах. Binz и Schulz~\cite{binz2023using} показали, что GPT-3 демонстрирует эффекты репрезентативности и предвзятости подтверждения на классических задачах из программы эвристик и предубеждений Канемана и Тверски. Suri и соавторы~\cite{suri2024large} расширили эту линию, исследовав эффект фрейминга и ошибку конъюнкции в GPT-4. Macmillan-Scott и Musolesi~\cite{macmillan2024irrationality} провели масштабный бенчмарк нескольких моделей на задаче Уэйсона. Тем не менее ни одна из указанных работ не рассматривает все три парадигмы~--- конъюнктивную ошибку, задачу выбора карточек и задачу открытия правила 2-4-6~--- в едином экспериментальном дизайне и не вводит непрерывных метрик, позволяющих измерить силу предубеждения за пределами бинарной классификации ответов.

\textbf{Цель и задачи исследования.}
Цель настоящей работы~--- систематически измерить и сравнить когнитивные предубеждения в современных LLM на трёх классических психологических парадигмах, охватывающих различные типы рассуждения: вероятностное суждение, условная логика и индуктивное открытие правил.

Для достижения данной цели поставлены следующие задачи:
\begin{itemize}
  \item провести анализ существующих подходов к изучению когнитивных предубеждений в LLM и выявить методологические ограничения предшествующих работ;
  \item разработать экспериментальный дизайн для трёх парадигм~--- задачи Линды (ошибка конъюнкции), задачи выбора карточек Уэйсона и задачи открытия правила 2-4-6~--- с контролем за контаминацией обучающих данных;
  \item сформировать оригинальные датасеты для каждого эксперимента, обеспечивающие структурное единообразие и отсутствие пересечения с обучающими корпусами;
  \item провести эксперименты на девяти современных языковых моделях, охватывающих широкий спектр архитектур, размеров и провайдеров;
  \item ввести и апробировать непрерывную метрику силы предубеждения (Severity of Bias) как альтернативу бинарной классификации ошибок;
  \item сравнить поведение моделей в трёх парадигмах и проверить гипотезу о масштабировании как факторе снижения предубеждений;
  \item на основании полученных результатов сформулировать рекомендации по выбору и оценке LLM для приложений, требующих нормативно корректного рассуждения.
\end{itemize}

\textbf{Объект и предмет исследования.}
Объектом исследования являются большие языковые модели как системы рассуждения. Предметом исследования служат когнитивные предубеждения в ответах LLM на задачи из трёх психологических парадигм: ошибка конъюнкции, предвзятость подтверждения в условной логике и предвзятость подтверждения в индуктивном открытии правил.

\textbf{Научная новизна.}
Научная новизна работы определяется следующими результатами:
\begin{itemize}
  \item впервые проведено сравнительное исследование трёх классических парадигм когнитивной психологии~--- ошибки конъюнкции, задачи Уэйсона и задачи 2-4-6~--- в едином экспериментальном дизайне на одном и том же наборе из девяти моделей, что позволяет рассматривать предубеждение как единый феномен на спектре типов рассуждения;
  \item введена оригинальная непрерывная метрика Severity of Bias, основанная на самооценке уверенности модели и позволяющая обнаруживать токен-предубеждение в случаях, когда бинарная метрика демонстрирует потолок точности;
  \item разработан оригинальный датасет из 130 биографических виньеток для задачи Линды с систематической факториальной подстановкой демографических переменных (имя, раса, возраст), что обеспечивает контроль за конфаундерами нарративного стиля;
  \item введено условие «фантастических социальных контрактов» в задаче Уэйсона, позволяющее диссоциировать эффект знакомости содержания от эффекта деонтической структуры, что, насколько известно автору, ранее не реализовывалось в исследованиях LLM;
  \item обнаружен и описан внутрисемейный парадокс масштаба в семействах моделей Gemma и Qwen: меньшая модель демонстрирует более низкий уровень предубеждения, чем большая модель того же семейства;
  \item разработан воспроизводимый экспериментальный фреймворк с унифицированным API-пайплайном через OpenRouter, оригинальными датасетами и статистическим протоколом с FDR-коррекцией, опубликованный в открытом доступе. Нетривиальность фреймворка определяется не объёмом реализации, а методологическими решениями на пересечении когнитивной психологии и ML: факториальным дизайном стимульного материала, формализацией оценки гипотез через экстенсиональную эквивалентность на основе Монте-Карло (IOU) и разработкой деконтаминированных экспериментальных условий~--- совокупность которых не воспроизводится без предметной экспертизы в обеих областях.
\end{itemize}

\textbf{Практическая значимость.}
Результаты исследования имеют непосредственное прикладное значение в нескольких направлениях. Во-первых, разработанный фреймворк тестирования когнитивных предубеждений может применяться при выборе языковой модели для систем поддержки принятия решений, где критична нормативная корректность рассуждения~--- в частности, в медицинской диагностике, юридическом анализе и финансовом консультировании. Во-вторых, введённая метрика Severity of Bias предоставляет более тонкий инструмент оценки, чем бинарная точность: например, \texttt{claude-sonnet-4.6} допускает лишь одну бинарную ошибку из 520 пар, однако метрика фиксирует delta\_bias $= +0{,}105$, что свидетельствует о латентном токен-предубеждении, невидимом при стандартном подходе. В-третьих, выявленные различия в поведении моделей на различных типах задач позволяют формулировать рекомендации по prompt engineering для снижения предубеждений: в частности, переход к reasoning-режиму обеспечивает качественный скачок в задаче 2-4-6 (CTR снижается ниже $0{,}63$), недостижимый средствами обычного промптинга. Полученные датасеты и экспериментальный код публикуются в открытом доступе в целях воспроизводимости исследований.

\textbf{Положения, выносимые на защиту.}
\begin{enumerate}
  \item Токен-предубеждение в форме ошибки конъюнкции является широко распространённым, но не абсолютно универсальным феноменом: оно обнаружено у восьми из девяти тестируемых моделей при использовании оригинального датасета из 130 биографических виньеток.
  \item Метрика Severity of Bias, основанная на самооценке уверенности модели, обнаруживает предубеждение там, где бинарная метрика точности достигает потолка, что подтверждено на примере \texttt{claude-sonnet-4.6} (delta\_bias $= +0{,}105$ при единственной бинарной ошибке из 520 пар).
  \item Предубеждение подтверждения в дедуктивном рассуждении воспроизводится у всех девяти моделей в нейтральном промпт-условии; деонтическая структура задачи (включая незнакомые фантастические правила) систематически снижает его интенсивность у восьми из девяти моделей.
  \item Предубеждение подтверждения в индуктивном открытии правил является универсальным для всех non-reasoning моделей (CTR $\geq 0{,}63$); переход к reasoning-режиму обеспечивает качественный скачок, недостижимый через prompt engineering.
  \item Масштаб модели не является монотонным предиктором устойчивости к когнитивным предубеждениям: в семействах Gemma и Qwen меньшая модель демонстрирует более низкий уровень предубеждения в задаче ошибки конъюнкции, чем большая.
  \item Профили когнитивных предубеждений по трём парадигмам не коррелируют монотонно на уровне отдельных моделей, что свидетельствует о мультидимензиональной природе данного феномена и невозможности его описания единой агрегированной метрикой.
\end{enumerate}

\textbf{Методологическая основа.}
Исследование опирается на классические парадигмы когнитивной психологии~\cite{tversky1983extensional, wason1966reasoning, wason1960failure}, адаптированные для тестирования языковых моделей. В качестве статистических инструментов используются точный тест МакНемара для согласованных пар~\cite{mcnemar1947note}, парный t-тест, критерий Уилкоксона и процедура Бенджамини-Хохберга для контроля частоты ложных открытий~\cite{benjamini1995controlling}. Для оценки качества гипотез в задаче 2-4-6 применяется метрика расширенного пересечения по объединению (IOU) на основе метода Монте-Карло (200\,000 троек), впервые предложенная Banatt и соавторами~\cite{banatt2024wilt} и расширенная в настоящей работе.

\textbf{Апробация результатов.}
Основные результаты исследования были представлены на кандидатском экзамене в AI Talent Hub ИТМО (январь 2026~г.). В ходе представления были получены замечания от научного руководителя и членов комиссии, часть которых учтена при доработке методологии эксперимента 3 и формулировке ограничений исследования. По результатам проведённых экспериментов подготовлена рукопись статьи, планируемая к подаче в рецензируемое издание в области искусственного интеллекта и когнитивных наук. Экспериментальный код и датасеты публикуются в открытом доступе для обеспечения воспроизводимости результатов.

\textbf{Личный вклад автора.}
Все этапы исследования выполнены автором самостоятельно: постановка задачи и разработка экспериментального дизайна, составление оригинальных датасетов (130 биографических виньеток с факториальными подстановками, 5 датасетов задачи Уэйсона по 100--101 записи, 50 правил для задачи 2-4-6), реализация экспериментального фреймворка на Python с асинхронным пайплайном через OpenRouter, проведение экспериментов на девяти языковых моделях, статистический анализ результатов, введение метрики Severity of Bias и интерпретация полученных данных.

\textbf{Использование систем искусственного интеллекта.}
При подготовке настоящей работы языковые модели использовались 
исключительно как вспомогательный инструмент на следующих этапах: 
редактирование и вычитка текстовых формулировок (проверка 
грамматики, стилистической связности и устранение повторов); 
выявление логических противоречий и несогласованностей между 
разделами; уточнение англоязычной терминологии и корректности 
перевода специализированных понятий; генерация текстового 
наполнения датасетов (биографических виньеток, формулировок 
правил) по авторским шаблонам, факториальному дизайну и 
заданным ограничениям через специализированные скрипты с 
последующей ручной верификацией каждого элемента; генерация 
шаблонного boilerplate-кода с последующей ручной проверкой 
и адаптацией.

Все содержательные решения~--- постановка задачи, выбор парадигм 
и моделей, разработка экспериментального и факториального дизайна, 
архитектура фреймворка, статистический анализ и интерпретация 
результатов~--- приняты и реализованы автором самостоятельно. 
Языковые модели не использовались для формулировки гипотез, 
анализа литературы или интерпретации экспериментальных данных.

\textbf{Структура работы.}
Работа состоит из введения, четырёх глав, заключения, списка использованных источников и приложений.

В \textit{первой главе} проводится анализ предметной области: рассматриваются теоретические основы когнитивных предубеждений, анализируется проблема их воспроизведения в LLM в соответствии с текущим состоянием области, выполняется обзор смежных работ и формулируются исследовательские гипотезы H1--H4.

Во \textit{второй главе} описывается дизайн экспериментов: обосновывается выбор парадигм, представлены оригинальные датасеты, промпт-условия, протестированные модели и система метрик; приводятся результаты экспериментов и их анализ.

В \textit{третьей главе} описывается архитектура воспроизводимого экспериментального фреймворка: структура кодовой базы, параметры запросов, механизм отказоустойчивости, реализация метрики IOU.

В \textit{четвёртой главе} проводится статистическая валидация гипотез H1--H4 с применением теста МакНемара и процедуры Бенджамини-Хохберга, обсуждаются ограничения исследования, приводится апробация результатов и намечаются направления дальнейших исследований.

В \textit{заключении} подводятся итоги работы и формулируются основные выводы.